\Article{13}

Factores membri secundi aequationis $[5]$ ita quoque exhiberi possunt

\[
\begin{aligned}
r-r^{-1}&=-(r^{n-1}-r^{-n+1})\\
r^3-r^{-3}&=-(r^{n-3}-r^{-n+3})\\
r^5-r^{-5}&=-(r^{n-5}-r^{-n+5}),\;\text{etc.}
\end{aligned}
\]

usque ad

\[
\begin{aligned}
r^{n-2}-r^{-n+2}=-(r^2-r^{-2})
\end{aligned}
\]

quo pacto aequatio ista hanc formam assumit:

\SourceLeaf{35}

\[
\begin{aligned}
W&=(-1)^{\frac12(n-1)}(r^2-r^{-2})(r^4-r^{-4})(r^6-r^{-6})\\&\qquad\dotsfour(r^{n-1}-r^{-n+1})
\end{aligned}
\]

Multiplicando hanc aequationem per $[5]$ in forma primitiva, prodit

\[
\begin{aligned}
W^2&=(-1)^{\frac12(n-1)}(r-r^{-1})(r^2-r^{-2})(r^3-r^{-3})\\&\qquad\dotsfour(r^{n-1}-r^{-n+1})
\end{aligned}
\]

ubi $(-1)^{\frac12(n-1)}$ est vel $=+1$ vel $=-1$, prout $n$ est formae $4\mu+1$, vel formae $4\mu+3$.\quad Hinc

\[
\begin{aligned}
W^2&=\pm r^{\frac12n(n-1)}(1-r^{-2})(1-r^{-4})(1-r^{-6})\\&\qquad\dotsfour(1-r^{-2(n-1)})
\end{aligned}
\]

Sed nullo negotio perspicitur, $r^{-2}$, $r^{-4}$, $r^{-6}\dotsfour r^{-2n+2}$ exhibere omnes radices aequationis $x^n-1=0$, radice $x=1$ excepta, unde locum habere debebit aequatio identica indefinita

\[
\begin{aligned}
&(x-r^{-2})(x-r^{-4})(x-r^{-6})\dotsfour(x-r^{-2n+2})\\&\qquad=x^{n-1}+x^{n-2}+x^{n-3}+\text{etc.}+x+1
\end{aligned}
\]

Quamobrem statuendo $x=1$, fiet

\[
\begin{aligned}
(1-r^{-2})(1-r^{-4})(1-r^{-6})\dotsfour(1-r^{-2n+2})=n
\end{aligned}
\]

et quum manifesto sit $r^{\frac12n(n-1)}=1$, aequatio nostra transit in hanc

\[
\begin{aligned}
W^2=\pm n\  .\  .\  .\  .\  .\  .\  .\  .\  .\  .\ [6]
\end{aligned}
\]

In casu itaque eo, ubi $n$ est formae $4\mu+1$, fiet

\[
\begin{aligned}
W=\pm\surd n,\quad\text{et proin}\quad T=\pm\surd n,\quad U=0
\end{aligned}
\]

Contra in casu altero, ubi $n$ est formae $4\mu+3$, fiet

\[
\begin{aligned}
W=\pm i\surd n,\quad\text{adeoque}\quad T=0,\quad U=\pm\surd n
\end{aligned}
\]

\Article{14}

Methodus art. praec. valorem tantummodo absolutum aggregatorum $T$, $U$ assignat, ambiguumque linquit, utrum statuere oporteat $T$ in casu priori atque $U$ in casu posteriori $=+\surd n$, an $=-\surd n$.\quad Hoc autem, saltem pro casu eo ubi $k=1$, ex aequatione $[5]$ sequenti modo decidere licebit.\quad Quum sit, pro $k=1$,

\SourceLeaf{36}

\[
\begin{aligned}
r-r^{-1}&=2i\sin\omega\\
r^3-r^{-3}&=2i\sin3\omega\\
r^5-r^{-5}&=2i\sin5\omega\;\text{etc.}
\end{aligned}
\]

aequatio ista transmutatur in

\[
\begin{aligned}
W=(2i)^{\frac12(n-1)}\sin\omega\sin3\omega\sin5\omega\ldots\sin(n-2)\omega
\end{aligned}
\]

Iam in casu eo, ubi $n$ est formae $4\mu+1$, in serie numerorum imparium

\[
\begin{aligned}
1,\;3,\;5,\;7\dotsfour\tfrac12(n-3),\quad\tfrac12(n+1)\dotsfour(n-2)
\end{aligned}
\]

reperiuntur $\frac14(n-1)$, qui sunt minores quam $\frac12n$, hisque manifesto respondent sinus positivi; contra reliqui $\frac14(n-1)$ erunt maiores quam $\frac12n$, hisque sinus negativi respondebunt: quapropter productum omnium sinuum statuendum est aequale producto e quantitate positiva in multiplicatorem $(-1)^{\frac14(n-1)}$, adeoque $W$ aequalis erit producto e quantitate reali positiva in $i^{n-1}$ sive in $1$, quoniam $i^4=1$, atque $n-1$ per $4$ divisibilis: i. e. quantitas $W$ erit realis positiva, unde necessario esse debebit

\[
\begin{aligned}
W=+\surd n,\quad T=+\surd n
\end{aligned}
\]

In casu altero, ubi $n$ est formae $4\mu+3$ in serie numerorum imparium

\[
\begin{aligned}
1,\;3,\;5,\;7\dotsfour\tfrac12(n-1),\quad\tfrac12(n+3)\dotsfour(n-2)
\end{aligned}
\]

priores $\frac14(n+1)$ erunt minores quam $\frac12n$, reliqui $\frac14(n-3)$ autem maiores.\quad Hinc inter sinus arcuum $\omega$, $3\omega$, $5\omega\dotsfour(n-2)\omega$ negativi erunt $\frac14(n-3)$, adeoque $W$ erit productum ex $i^{\frac12(n-1)}$ in quantitatem realem positivam in $(-1)^{\frac14(n-3)}$; factor tertius est $=i^{\frac12(n-3)}$, qui cum primo iunctus producit $i^{n-2}=i$, quoniam $i^{n-3}=1$.\quad Quamobrem necessario erit

\[
\begin{aligned}
W=+i\surd n,\quad\text{atque}\quad U=+\surd n
\end{aligned}
\]

\Article{15}

Iam ostendemus, quo pacto eaedem conclusiones e progressione in art. $9$ considerata deduci possint.\quad Scribamus in aequ. $[4]$ pro $x^{\frac12}$, $-y^{-1}$, eritque

\SourceLeaf{37}

\[
\begin{aligned}
&1-y^{-1}\frac{1-y^{-2m}}{1-y^{-2}}+y^{-2}\frac{(1-y^{-2m})(1-y^{-2m+2})}{(1-y^{-2})(1-y^{-4})}\\&\qquad-y^{-3}\frac{(1-y^{-2m})(1-y^{-2m+2})(1-y^{-2m+4})}{(1-y^{-2})(1-y^{-4})(1-y^{-6})}+\text{etc.}
\end{aligned}
\]

usque ad terminum $m+1^{\mathrm{tum}}$

\[
\begin{aligned}
&=(1-y^{-1})(1+y^{-2})(1-y^{-3})(1+y^{-4})\\&\qquad\ldots(1\pm y^{-m})\  .\  .\  .\  .\ [7]
\end{aligned}
\]

Quodsi hic pro $y$ accipitur radix propria aequationis $y^n-1=0$, puta $r$, atque simul statuitur $m=n-1$, erit

\[
\begin{aligned}
\frac{1-y^{-2m}}{1-y^{-2}}&=\frac{1-r^2}{1-r^{-2}}=-r^2\\
\frac{1-y^{-2m+2}}{1-y^{-4}}&=\frac{1-r^4}{1-r^{-4}}=-r^4\\
\frac{1-y^{-2m+4}}{1-y^{-6}}&=\frac{1-r^6}{1-r^{-6}}=-r^6
\end{aligned}
\]

usque ad

\[
\begin{aligned}
\frac{1-y^{-2}}{1-y^{-2m}}=\frac{1-r^{2n-2}}{1-r^{-2n+2}}=-r^{2n-2}
\end{aligned}
\]

ubi notandum, nullum denominatorum $1-r^{-2}$, $1-r^{-4}$ etc. fieri $=0$.\quad Hinc aequatio $[7]$ hancce formam assumit

\[
\begin{aligned}
&1+r+r^4+r^9+\text{etc.}+r^{(n-1)^2}\\&\qquad=(1-r^{-1})(1+r^{-2})(1-r^{-3})\dotsfour(1+r^{-n+1})
\end{aligned}
\]

Multiplicando in membro secundo huius aequationis terminum primum per ultimum, secundum per penultimum etc., habemus

\[
\begin{aligned}
(1-r^{-1})(1+r^{-n+1})&=r-r^{-1}\\
(1+r^{-2})(1-r^{-n+2})&=r^{n-2}-r^{-n+2}\\
(1-r^{-3})(1+r^{-n+3})&=r^3-r^{-3}\\
(1+r^{-4})(1-r^{-n+4})&=r^{n-4}-r^{-n+4}\;\text{etc.}
\end{aligned}
\]

Ex his productis partialibus facile perspicietur conflari productum

\[
\begin{aligned}
&(r-r^{-1})(r^3-r^{-3})(r^5-r^{-5})\\&\qquad\dotsfour(r^{n-4}-r^{-n+4})(r^{n-2}-r^{-n+2})
\end{aligned}
\]

quod itaque erit

\[
\begin{aligned}
=1+r+r^4+r^9+\text{etc.}+r^{(n-1)^2}=W
\end{aligned}
\]

Haec aequatio identica est cum aequ. $[5]$ in art. $12$ e progressione prima derivata, ratiociniaque dein reliqua eodem modo adstruentur, ut in artt. $13$ et $14$.

\SourceLeaf{38}

\Article{16}

Transimus ad casum alterum, ubi $n$ est numerus par.\quad Sit primo $n$ formae $4\mu+2$ sive impariter par, patetque, numeros $\frac14nn$, $(\frac12n+1)^2-1$, $(\frac12n+2)^2-4$ etc. sive generaliter $(\frac12n+\lambda)^2-\lambda\lambda$ per $\frac12n$ divisos producere quotientes impares, adeoque secundum modulum $n$ congruos fieri ipsi $\frac12n$.\quad Hinc colligitur, si $r$ sit radix propria aequationis $x^n-1=0$, adeoque $r^{\frac12n}=-1$, fieri

\[
\begin{aligned}
r^{(\frac12n)^2}&=-1\\
r^{(\frac12n+1)^2}&=-r\\
r^{(\frac12n+2)^2}&=-r^4\\
r^{(\frac12n+2)^2}&=-r^9\;\text{etc.}
\end{aligned}
\]

Hinc in progressione

\[
\begin{aligned}
1+r+r^4+r^9+\text{etc.}+r^{(n-1)^2}
\end{aligned}
\]

terminus $r^{(\frac12n)^2}$ destruet primum, sequens secundum etc., adeoque erit

\[
\begin{aligned}
W=0,\quad T=0,\quad U=0
\end{aligned}
\]

\Article{17}

Superest casus, ubi $n$ est formae $4\mu$ sive pariter par.\quad Hic generaliter $(\frac12n+\lambda)^2-\lambda\lambda$ divisibilis erit per $n$, adeoque

\[
\begin{aligned}
r^{(\frac12n+\lambda)^2}=r^{\lambda\lambda}
\end{aligned}
\]

Hinc in serie

\[
\begin{aligned}
1+r+r^4+r^9+\text{etc.}+r^{(n-1)^2}
\end{aligned}
\]

terminus $r^{(\frac12n)^2}$ aequalis erit primo, sequens secundo etc., ita ut fiat

\[
\begin{aligned}
W=2(1+r+r^4+r^9+\text{etc.}+r^{(\frac12n-1)^2})
\end{aligned}
\]

Iam supponamus, in aequ. $[7]$ art. 15 statui $m=\frac12n-1$, et pro $y$ accipi radicem propriam aequationis $y^n-1=0$, puta $r$.\quad Tunc perinde ut in art. 15 aequatio sequentem formam obtinet:

\[
\begin{aligned}
&1+r+r^4+\text{etc.}+r^{(\frac12n-1)^2}\\&\qquad=(1-r^{-1})(1+r^{-2})(1-r^{-3})\dotsfour(1-r^{-\frac12n+1})
\end{aligned}
\]

\SourceLeaf{39}

sive

\[
\begin{aligned}
W&=2(1-r^{-1})(1+r^{-2})(1-r^{-3})(1+r^{-4})\\&\qquad\ldots(1-r^{-\frac12n+1})\quad\dotstwo\quad[8]
\end{aligned}
\]

Porro quum sit $r^{\frac12n}=-1$, adeoque

\[
\begin{aligned}
1+r^{-2}&=-r^{\frac12n-2}(1-r^{-\frac12n+2})\\
1+r^{-4}&=-r^{\frac12n-4}(1-r^{-\frac12n+4})\\
1+r^{-6}&=-r^{\frac12n-6}(1-r^{-\frac12n+6})\;\text{etc.}
\end{aligned}
\]

productumque e factoribus $-r^{\frac12n-2}$, $-r^{\frac12n-4}$, $-r^{\frac12n-6}$ etc. usque ad $-r^2$ fiat $=(-1)^{\frac14n-1}\,r^{\frac1{16}nn-\frac14n}$, aequatio praecedens ita quoque exhiberi potest

\[
\begin{aligned}
W&=2(-1)^{\frac14n-1}\,r^{\frac1{16}nn-\frac14n}(1-r^{-1})(1-r^{-2})\\&\qquad(1-r^{-3})(1-r^{-4})\dotsfour(1-r^{-\frac12n+1})
\end{aligned}
\]

Quum habeatur

\[
\begin{aligned}
1-r^{-1}&=-r^{-1}(1-r^{-n+1})\\
1-r^{-2}&=-r^{-2}(1-r^{-n+2})\\
1-r^{-3}&=-r^{-3}(1-r^{-n+3})\;\text{etc.}
\end{aligned}
\]

erit

\[
\begin{aligned}
&(1-r^{-1})(1-r^{-2})(1-r^{-3})\ldots(1-r^{-\frac12n+1})\\
&\quad=(-1)^{\frac12n-1}\,r^{-\frac18nn+\frac14n}(1-r^{-\frac12n-1})\\&\qquad(1-r^{-\frac12n-2})(1-r^{-\frac12n-3})\ldots(1-r^{-n+1})
\end{aligned}
\]

adeoque

\[
\begin{aligned}
W&=2(-1)^{\frac34n-2}\,r^{-\frac1{16}nn}(1-r^{-\frac12n-1})\\&\qquad(1-r^{-\frac12n-2})(1-r^{-\frac12n-3})\dotsfour(1-r^{-n+1})
\end{aligned}
\]

Multiplicando hunc valorem ipsius $W$ per prius inventum, adiungendoque utrimque factorem $1-r^{-\frac12n}$, prodit

\[
\begin{aligned}
(1-r^{-\frac12n})W^2&=4(-1)^{n-3}r^{-\frac14n}(1-r^{-1}\mathord{.})\\&\qquad(1-r^{-2})(1-r^{-3})\dotsfour(1-r^{-n+1})
\end{aligned}
\]

Sed fit

\[
\begin{aligned}
1-r^{-\frac12n}&=2\\
(-1)^{n-3}&=-1\\
r^{-\frac14n}&=-r^{\frac14n}
\end{aligned}
\]

\[
\begin{aligned}
(1-r^{-1})(1-r^{-2})(1-r^{-3})\dotsfour(1-r^{-n+1})=n
\end{aligned}
\]

Unde tandem concluditur

\SourceLeaf{40}

\[
\begin{aligned}
W^2=2r^{\frac14n}n\  .\  .\  .\  .\  .\  .\  .\  .\  .\ [9]
\end{aligned}
\]

Iam facile perspicietur, $r^{\frac14n}$ esse vel $=+i$ vel $=-i$, prout scilicet $k$ vel formae $4\mu+1$ sit, vel formae $4\mu+3$.\quad Et quum sit

\[
\begin{aligned}
2i=(1+i)^2,\quad -2i=(1-i)^2
\end{aligned}
\]

erit in casu eo, ubi $k$ est formae $4\mu+1$,

\[
\begin{aligned}
W=\pm(1+i)\surd n,\quad\text{adeoque}\quad T=U=\pm\surd n
\end{aligned}
\]

in casu altero autem, ubi $k$ est formae $4\mu+3$

\[
\begin{aligned}
W=\pm(1-i)\surd n,\quad\text{adeoque}\quad T=-U=\pm\surd n
\end{aligned}
\]
